\documentclass[UTF8,11pt,a4paper]{ctexart}
\usepackage[margin=24mm]{geometry}
\usepackage{amsmath,amssymb,amsthm,mathtools,bm}
\usepackage{tikz-cd,enumitem,longtable,booktabs}
\usepackage{xurl}
\usepackage[unicode,hidelinks]{hyperref}
\allowdisplaybreaks
\setlength{\emergencystretch}{3em}
\setlength{\parindent}{0pt}
\setlength{\parskip}{0.45em}
\newtheorem*{theorem}{Theorem}
\newtheorem*{lemma}{Lemma}
\newtheorem*{proposition}{Proposition}
\newtheorem*{corollary}{Corollary}
\theoremstyle{definition}\newtheorem*{definition}{Definition}
\theoremstyle{remark}\newtheorem*{remark}{Remark}
\DeclareMathOperator{\Hom}{Hom}
\DeclareMathOperator{\End}{End}
\DeclareMathOperator{\Tr}{Tr}
\DeclareMathOperator{\rank}{rank}
\DeclareMathOperator{\Ind}{Ind}
\DeclareMathOperator{\Res}{Res}
\DeclareMathOperator{\im}{Im}
\DeclareMathOperator{\id}{id}
\newcommand{\R}{\mathbb R}
\newcommand{\C}{\mathbb C}
\newcommand{\Q}{\mathbb Q}
\newcommand{\Z}{\mathbb Z}
\newcommand{\N}{\mathbb N}
\newcommand{\E}{\mathbb E}
\newcommand{\Prob}{\mathbb P}
\newcommand{\eps}{\varepsilon}
\begin{document}
\section*{095｜按置换循环数定义的群行列式的整根分解}
\textbf{作者／集体来源：}Darij Grinberg

\textbf{作品：}A representation-theoretical solution to MathOverflow question \#88399, version 3 (4 December 2013; minor corrections 9 June 2026)

\textbf{原文入口：}\url{https://www.cip.ifi.lmu.de/~grinberg/algebra/MO88399.pdf}

\textbf{原文位置与计题范围：}Theorems 1--2及§2--§3完整论证（页2--8），含脚注3--6中的因子分解、首一性和张量置换迹计算。只计Theorem 1这一行列式问题。

\textbf{获取日期：}2026-09-28。\quad\textbf{复用依据：}CC0-1.0，作者作品页对应条目明确注明。

\textbf{整理归属：}下列题面与证明取自所列人类来源，按注明范围转录、摘编；LaTeX环境、标题、换行及中文说明由助手整理。编辑调整在题内说明。

\section*{作者命题与人类证明}

\subsection*{Statement of the problem (\S1)}
Let $n\in\N$. Consider the $n$th symmetric group $S_n$. For every $w\in S_n$, let $\sigma(w)$ denote the number of cycles in the cycle decomposition of the permutation $w$ (this includes cycles consisting of one element). We can consider the matrix $(x^{\sigma(gh^{-1})})_{g,h\in S_n}$; this is a matrix over the polynomial ring $\Q[x]$, whose rows and whose columns are indexed by the elements of $S_n$.
\begin{theorem}[1]
The polynomial
\[
\det\bigl((x^{\sigma(gh^{-1})})_{g,h\in S_n}\bigr)\in\Q[x]
\]
factors into linear factors of the form $x-\ell$ with $\ell\in\{-n+1,-n+2,\ldots,n-1\}$.
\end{theorem}

\subsection*{Reduction to representation theory (\S2)}
For any finite group $G$, let $\operatorname{Irrep}G$ denote a set of representatives of all irreducible representations of $G$ over $\C$ modulo isomorphism. From the theory of group determinants (the results of [1], or the proof of Theorem 4.7 in [2]), we know that if $G$ is a finite group, and $X_g$ is an indeterminate for every $g\in G$, then
\[
\det\bigl((X_{gh^{-1}})_{g,h\in G}\bigr)
=\prod_{\rho\in\operatorname{Irrep}G}
\left(\det\left(\sum_{g\in G}\rho(g)X_g\right)\right)^{\dim\rho}.
\]
Applying this to $G=S_n$ and evaluating this polynomial identity at $X_g=x^{\sigma(g)}$, we obtain
\begin{equation}\label{gr:groupdet}
\det\bigl((x^{\sigma(gh^{-1})})_{g,h\in S_n}\bigr)
=\prod_{\rho\in\operatorname{Irrep}S_n}
\left(\det\left(\sum_{g\in S_n}\rho(g)x^{\sigma(g)}\right)\right)^{\dim\rho}.
\end{equation}
Hence, in order to show that the polynomial on the left factors as claimed, it is enough to prove that, for every irreducible representation $\rho$ of $S_n$ over $\C$, the polynomial $\det(\sum_{g\in S_n}\rho(g)x^{\sigma(g)})$ factors into such linear factors. We are going to show something better.
\begin{theorem}[2]
Let $\lambda=(\lambda_1,\lambda_2,\ldots,\lambda_n)$ be a partition of $n$. Let $m_\lambda$ be the number of nonzero parts of $\lambda$. Let $\rho_\lambda$ be the irreducible representation of $S_n$ over $\C$ corresponding to $\lambda$. Then
\begin{equation}\label{gr:scalar}
\sum_{g\in S_n}\rho_\lambda(g)x^{\sigma(g)}
=\frac{n!}{\dim\rho_\lambda}
\prod_{1\le i<j\le m_\lambda}\frac{\lambda_i-\lambda_j+j-i}{j-i}
\prod_{i=1}^{m_\lambda}
\frac{\binom{x+\lambda_i-i}{\lambda_i}}
{\binom{\lambda_i+m_\lambda-i}{m_\lambda-i}}
\cdot\id_{\rho_\lambda}.
\end{equation}
\end{theorem}
\begin{proof}[Proof of Theorem 1 using Theorem 2]
For every partition $\lambda$ of $n$, denote by $\rho_\lambda$ the irreducible representation corresponding to $\lambda$, and by $m_\lambda$ the number of its nonzero parts. The isomorphism classes of irreducible representations of $S_n$ over $\C$ are in one-to-one correspondence with the partitions of $n$.

For the following displayed calculation only, abbreviate the scalar on the right of \eqref{gr:scalar} by $c_\lambda(x)$ (editorial abbreviation). Then
\begin{align*}
\prod_{\rho\in\operatorname{Irrep}S_n}
\left(\det\sum_{g\in S_n}\rho(g)x^{\sigma(g)}\right)^{\dim\rho}
&=\prod_{\lambda\vdash n}
\left(\det\sum_{g\in S_n}\rho_\lambda(g)x^{\sigma(g)}\right)^{\dim\rho_\lambda}\\
&=\prod_{\lambda\vdash n}\bigl(\det(c_\lambda(x)\id_{\rho_\lambda})\bigr)^{\dim\rho_\lambda}\\
&=\prod_{\lambda\vdash n}\bigl(c_\lambda(x)^{\dim\rho_\lambda}\bigr)^{\dim\rho_\lambda}
=\prod_{\lambda\vdash n}c_\lambda(x)^{(\dim\rho_\lambda)^2}.
\end{align*}
Here we used \eqref{gr:scalar} and $\det(z\id_d)=z^d$. Combined with \eqref{gr:groupdet}, this yields
\begin{equation}\label{gr:detanswer}
\det\bigl((x^{\sigma(gh^{-1})})_{g,h\in S_n}\bigr)
=\prod_{\lambda\vdash n}
\left(
\frac{n!}{\dim\rho_\lambda}
\prod_{1\le i<j\le m_\lambda}\frac{\lambda_i-\lambda_j+j-i}{j-i}
\prod_{i=1}^{m_\lambda}
\frac{\binom{x+\lambda_i-i}{\lambda_i}}
{\binom{\lambda_i+m_\lambda-i}{m_\lambda-i}}
\right)^{(\dim\rho_\lambda)^2}.
\end{equation}
The right hand side is clearly a polynomial in $x$ which factors into a product of a constant and linear factors. All of the linear factors have the form $x+\lambda_i-i-\alpha$ for $\alpha\in\{0,1,\ldots,\lambda_i-1\}$ for various partitions $\lambda$ of $n$ and various $i\in\{1,\ldots,m_\lambda\}$: the only place where $x$ occurs is
\[
\binom{x+\lambda_i-i}{\lambda_i}
=\frac{(x+\lambda_i-i)(x+\lambda_i-i-1)\cdots(x+\lambda_i-i-(\lambda_i-1))}{\lambda_i!}.
\]
By very simple combinatorics, each of these factors has the form $x-\ell$ for $\ell\in\{-n+1,-n+2,\ldots,n-1\}$. Moreover, the constant is $1$ because the determinant polynomial is monic. To see this, write the determinant as a sum over permutations of the set $S_n$ (not permutations in $S_n$): the highest degree of $x$ is contributed by the product of the main diagonal. The main diagonal is filled with $x^{\sigma(\id)}=x^n$, while all other entries are lower powers of $x$. Thus Theorem 1 is proven using Theorem 2.
\end{proof}

\subsection*{Proof of Theorem 2 (\S3)}
\begin{proof}
First of all, \eqref{gr:scalar} is a polynomial identity in $x$. Hence, we can WLOG assume that $x$ is not a polynomial indeterminate in $\Q[x]$, but an integer greater than $n$ (because if a polynomial identity over $\Q$ holds for infinitely many integers, then it must always hold). Assume this. Since $x$ is an integer greater than $n$, this allows us to find a $\C$-vector space of dimension $x$. Let $V$ be such a vector space.

For every $S_n$-module $P$, let $\chi_P$ denote its character. Note that every $h\in S_n$ satisfies
\begin{equation}\label{gr:cycletrace}
\chi_{V^{\otimes n}}(h)=x^{\sigma(h)}.
\end{equation}
\paragraph{Proof of \eqref{gr:cycletrace} (source's footnote 5).}
Denote the action of $h$ on $V^{\otimes n}$ by $h|_{V^{\otimes n}}$. By the definition of a character, $\chi_{V^{\otimes n}}(h)=\Tr(h|_{V^{\otimes n}})$.
Pick a basis $(e_1,e_2,\ldots,e_x)$ of $V$. This induces the basis
\[
(e_{i_1}\otimes e_{i_2}\otimes\cdots\otimes e_{i_n})_{(i_1,\ldots,i_n)\in\{1,\ldots,x\}^n}
\]
of $V^{\otimes n}$. By the definition of the action,
\[
h(e_{i_1}\otimes\cdots\otimes e_{i_n})
=e_{i_{h^{-1}(1)}}\otimes\cdots\otimes e_{i_{h^{-1}(n)}}.
\]
Thus the linear map is represented by the permutation matrix of the permutation $h^{(\times n)}$ which sends $(i_1,\ldots,i_n)$ to $(i_{h^{-1}(1)},\ldots,i_{h^{-1}(n)})$. Its trace is the number of fixed points of $h^{(\times n)}$.
An $n$-tuple is fixed if and only if every $j$ satisfies $i_j=i_{h^{-1}(j)}$, or equivalently, each pair $j,k$ in the same cycle of $h$ satisfies $i_j=i_k$. To choose a fixed point, we need only specify one element of $\{1,\ldots,x\}$ for each cycle of $h$; these choices are arbitrary and independent, and besides them there is no more freedom. Thus there are precisely $x^{\sigma(h)}$ fixed points. This proves \eqref{gr:cycletrace}.

Let $L_\lambda$ be the representation of $\operatorname{GL}(V)$ corresponding to $\lambda$, that is, the image of $V$ under the $\lambda$th Schur functor. By one of the definitions of Schur functors,
$L_\lambda=\Hom_{\C[S_n]}(\rho_\lambda,V^{\otimes n})$, so
\begin{align}
\dim L_\lambda
&=\dim\Hom_{\C[S_n]}(\rho_\lambda,V^{\otimes n})
=\langle\chi_{V^{\otimes n}},\chi_{\rho_\lambda}\rangle\notag\\
&=\frac1{n!}\sum_{g\in S_n}\chi_{\rho_\lambda}(g)\chi_{V^{\otimes n}}(g^{-1})\notag\\
&=\frac1{n!}\sum_{g\in S_n}\Tr(\rho_\lambda(g))x^{\sigma(g^{-1})}\notag\\
&=\frac1{n!}\Tr\left(\sum_{g\in S_n}\rho_\lambda(g)x^{\sigma(g^{-1})}\right)
=\frac1{n!}\Tr\left(\sum_{g\in S_n}\rho_\lambda(g)x^{\sigma(g)}\right).\label{gr:trace}
\end{align}
We used the character inner product (Theorem 3.8 of [2]), \eqref{gr:cycletrace}, and $\sigma(g^{-1})=\sigma(g)$.
On the other hand, Theorem 4.63 of [2] (the Weyl character formula) yields, setting $\lambda_\ell=0$ for $\ell>m_\lambda$,
\begin{align*}
\dim L_\lambda
&=\prod_{1\le i<j\le x}\frac{\lambda_i-\lambda_j+j-i}{j-i}\\
&=\prod_{1\le i<j\le m_\lambda}\frac{\lambda_i-\lambda_j+j-i}{j-i}
\prod_{1\le i\le m_\lambda<j\le x}\frac{\lambda_i+j-i}{j-i}
\prod_{m_\lambda<i<j\le x}1\\
&=\prod_{1\le i<j\le m_\lambda}\frac{\lambda_i-\lambda_j+j-i}{j-i}
\prod_{i=1}^{m_\lambda}\prod_{j=m_\lambda+1}^x\frac{\lambda_i+j-i}{j-i}\\
&=\prod_{1\le i<j\le m_\lambda}\frac{\lambda_i-\lambda_j+j-i}{j-i}
\prod_{i=1}^{m_\lambda}
\frac{\binom{x+\lambda_i-i}{\lambda_i}}
{\binom{\lambda_i+m_\lambda-i}{m_\lambda-i}}.
\end{align*}
The last finite-product identity is straightforward to check. Combined with \eqref{gr:trace}, this gives
\begin{equation}\label{gr:scalartrace}
\Tr\left(\sum_{g\in S_n}\rho_\lambda(g)x^{\sigma(g)}\right)
=n!\prod_{1\le i<j\le m_\lambda}\frac{\lambda_i-\lambda_j+j-i}{j-i}
\prod_{i=1}^{m_\lambda}
\frac{\binom{x+\lambda_i-i}{\lambda_i}}
{\binom{\lambda_i+m_\lambda-i}{m_\lambda-i}}.
\end{equation}
But $\sum_{g\in S_n}gx^{\sigma(g)}$ is a central element of $\C[S_n]$, since $g\mapsto\sigma(g)$ is a class function (conjugate permutations have the same number of cycles). It therefore acts on any irreducible representation of $S_n$ as a scalar multiple of the identity, by Schur's lemma. In particular,
\[
\sum_{g\in S_n}\rho_\lambda(g)x^{\sigma(g)}
=\rho_\lambda\left(\sum_{g\in S_n}gx^{\sigma(g)}\right)
=\kappa\id_{\rho_\lambda}
\]
for some $\kappa\in\C$. Taking traces gives
\[
\Tr\left(\sum_{g\in S_n}\rho_\lambda(g)x^{\sigma(g)}\right)=\kappa\dim\rho_\lambda.
\]
Combining with \eqref{gr:scalartrace}, we find
\[
\kappa=\frac{n!}{\dim\rho_\lambda}
\prod_{1\le i<j\le m_\lambda}\frac{\lambda_i-\lambda_j+j-i}{j-i}
\prod_{i=1}^{m_\lambda}
\frac{\binom{x+\lambda_i-i}{\lambda_i}}
{\binom{\lambda_i+m_\lambda-i}{m_\lambda-i}}.
\]
Substitution yields \eqref{gr:scalar}. This proves Theorem 2.
\end{proof}
\paragraph{References retained from the source.}
[1] Keith Conrad, \emph{The Origin of Representation Theory}.
[2] Pavel Etingof et al., \emph{Introduction to Representation Theory}, arXiv:0901.0827v5.

\section*{整理说明与可修改标注（助手）}
\textbf{标准先修：}群行列式的一般分解；对称群不可约表示的分拆参数化；Schur函子与Weyl维数公式；Schur引理。

\textbf{本题仍需完成的工作：}将循环数权函数实现为张量幂特征标，通过Schur函子的维数确定中心元标量，再回到大行列式并定位所有整根。

\textbf{取材说明：}原例子和练习未计入；核心脚注已转入正文。为避免反复超长公式，Theorem 1的代入计算临时采用明确标出的标量缩写；所有等式步骤及维数的平方这一重数保留。文字作标点和重复措辞整理，不另造证明。原PDF文字层完整取得，网页截图接口未返回图像；二项式分式依文字层的完整上下结构转录。

\textbf{$c$：}组合矩阵的行列式分解

\textbf{$a$：}群代数中心元；群行列式；张量置换的迹；Schur函子；Weyl维数公式；多项式恒等式

\texttt{cross\_theory\_status = yes}

\textbf{具体联系及位置：}§3用整数x作为向量空间维数，将x的循环数次幂改写成张量幂迹；由另一群GL(V)的表示维数读回对称群中心元的作用，从而得到原组合矩阵的线性因子。

这些标注是非穷尽初稿；未列出不表示未使用。
\end{document}
